% !PRINTED 646
\weberpage{646}
%<<CONT
% [sic] Print changes the argument list here to $\xi_0,\ldots,\xi_{n-2}$;
% the function introduced on printed p. 645 has arguments through $\xi_{n-1}$.
2. ändert sich diese Function aber auch nicht, wenn unter den $k_\nu$ die cyklische Permutation $(k_0, k_1 \ldots k_{n-2})$ vorgenommen wird, und weil nun die Grössen $k_\nu$ die Wurzeln einer cyklischen Gleichung im Körper $\mathfrak{K}$ sind, so ist die Function $S(\xi_0, \xi_1 \ldots \xi_{n-2})$ eine Grösse in $\mathfrak{K}$, d. h. rational. Wendet man dies auf die Coëfficienten der Gleichung an, deren Wurzeln die Grössen (10) sind, so folgt, dass diese Grössen die Wurzeln einer Gleichung $n^{\mathrm{ten}}$ Grades in $\mathfrak{K}$ sind.

Was die Irreducibilität dieser Gleichung betrifft, so ist darüber Folgendes zu bemerken:

Die Radicale $\sqrt[n]{R_\nu}$ sind, wenn eines von ihnen rational ist, alle rational, und wenn eines irrational ist, sind es die anderen auch [§.\,185, (8)]. Welcher der beiden Fälle eintritt, hängt also nach §.\,179 davon ab, ob die Function $x^n - R_0$ in $\mathfrak{K}$ reducibel oder irreducibel ist. Ist sie reducibel, so ist eine der Grössen $\xi_h$ im Körper $\mathfrak{K}$ rational, und die Gleichung für $\xi$ ist reducibel. Alle Grössen $\xi$ sind in dem Körper $\mathfrak{K}(\varepsilon)$ enthalten, dessen Grad höchstens gleich $n-1$ ist.

Ist aber $\sqrt[n]{R_0}$ irrational, so ist $x^n - R_0$ in $\mathfrak{K}$ irreducibel. Jede Gleichung $\Phi(\xi_0) = 0$ geht durch Substitution des Werthes von $\xi_0$ in eine Gleichung $\Psi\!\left(\sqrt[n]{R_0}\right) = 0$ über, deren Coëfficienten in $\mathfrak{K}$ enthalten sind, und die also bestehen bleiben muss, wenn $\sqrt[n]{R_0}$ durch irgend einen seiner $n$ Werthe ersetzt wird.

Dadurch geht aber $\xi_0$ in jeden der Werthe $\xi_0, \xi_1 \ldots \xi_{n-1}$ über, d. h. die Gleichung $n^{\mathrm{ten}}$ Grades, deren Wurzeln die $n$ Grössen $\xi$ sind, ist in $\mathfrak{K}$ irreducibel.

Wir können also das Theorem I. so umkehren:

\gesperrt{II. Jede in der Form (9) enthaltene Grösse $\xi$ ist die Wurzel einer Gleichung $n^{\mathrm{ten}}$ Grades in $\mathfrak{K}$, und diese Gleichung ist irreducibel, wenn man von dem speciellen Falle absieht, dass eine ihrer Wurzeln rational, die anderen in $\mathfrak{K}(\varepsilon)$ enthalten sind.}
